\chapter{Nokta Tahmini ve Tahmin Edicinin Niteliği}
\index{nokta tahmini}
\index{tahmin edici}

\begin{hedefler}
Bu bölümün sonunda okuyucu:
\begin{itemize}
  \item araştırma sorusuna karşılık gelen hedef parametreyi açıkça tanımlayabilecek,
  \item parametre, tahmin edici, tahmin ve tahmin hatasını ayırt edebilecek,
  \item ortalama, oran, varyans, fark ve korelasyon için temel nokta tahminlerini hesaplayabilecek,
  \item yanlılık, varyans, standart hata, MSE, tutarlılık ve sağlamlık ölçütlerini yorumlayabilecek,
  \item örnekleme tasarımının ağırlıklı tahmin üzerindeki etkisini açıklayabilecek,
  \item nokta tahminini belirsizliği ve PDR bağlamıyla birlikte raporlayabilecektir.
\end{itemize}
\end{hedefler}

\begin{hazirlik}
Aynı evren için iki örneklemden 68 ve 74 ortalamaları elde edilmesi bir veri
hatası mıdır? Örnekleme değişkenliğiyle ilişkilendirin.
\end{hazirlik}

\section{PDR vakası: Yalnızlık düzeyi ve yardım gereksinimini tahmin etmek}

Bir üniversitenin psikolojik danışma merkezi iki bilgiye gereksinim duymaktadır:

\begin{enumerate}
  \item Birinci sınıf öğrencilerinin ortalama yalnızlık puanı nedir?
  \item Yüksek risk eşiğinin üzerinde puan alan öğrencilerin oranı nedir?
\end{enumerate}

Bütün öğrencileri ölçmek mümkün olmadığı için olasılıklı yöntemle 120 öğrenci
seçilmiştir. Örneklem ortalaması 43,6, standart sapması 9,8; yüksek risk
grubundaki öğrenci sayısı 36 bulunmuştur. Buna göre ortalama için nokta tahmini
$\bar{x}=43{,}6$, yüksek risk oranı için
\[
\hat p=\frac{36}{120}=0{,}30
\]
olur.

Bu iki sayı örneklemde gözlenen sonucu özetler; evren değerlerini kesin olarak
bildirmez. Başka bir 120 kişilik örneklem biraz farklı ortalama ve oran
verecektir. Tahminin niteliği, kullanılan formül kadar örnekleme ve ölçme
sürecine de bağlıdır.

\begin{researchbox}
``Yüksek risk oranı'' ancak risk eşiği, ölçme aracının geçerliği, hedef evren
ve ölçüm zamanı açıkça tanımlandığında anlamlı bir parametredir. İstatistiksel
tahmin, klinik sınıflandırmanın niteliğini tek başına güvenceye almaz.
\end{researchbox}

\section{Önce tahmin hedefini tanımlamak}
\index{hedef büyüklük}
\index{estimand@\textit{estimand}}

Tahmin edilmek istenen evren özelliğine \emph{hedef büyüklük} veya
\emph{estimand} denir. ``Öğrencilerin yalnızlık düzeyi'' ifadesi tek başına
yeterince açık değildir. Hedef şu bileşenleri belirtmelidir
\cite{casella2002,hogg2019}:

\begin{itemize}
  \item hangi öğrenci evreni,
  \item hangi ölçme aracı ve puanlama yöntemi,
  \item hangi zaman dönemi,
  \item ortalama, medyan, oran veya fark gibi hangi özet,
  \item eksik ve kapsam dışı gözlemlerin nasıl ele alınacağı.
\end{itemize}

Örneğin ``2026--2027 güz döneminde kayıtlı birinci sınıf öğrencilerinin X
Yalnızlık Ölçeği toplam puanı ortalaması'' açık bir hedefe daha yakındır.
``Programa katılanların son-test ortalaması'' ile ``programa atanmış bütün
öğrencilerin son-test ortalaması'' farklı hedefler olabilir.

\section{Parametre, tahmin edici ve tahmin}
\index{parametre}
\index{tahmin edici}
\index{tahmin}
\index{tahmin hatası}

Parametre evrene ait sabit fakat bilinmeyen bir niceliktir. Tahmin edici,
örneklemden hesaplanan rassal bir kuraldır; tahmin ise gözlenen örneklem için
bu kuralın aldığı sayısal değerdir. Örneğin $\bar X$ evren ortalaması $\mu$
için bir tahmin edici, $\bar x=72{,}4$ ise elde edilen tahmindir.

Tahmin hatası
\[
\hat\theta-\theta
\]
ile gösterilir. Gerçek parametre çoğu uygulamada bilinmediği için tek bir
araştırmada bu hata doğrudan hesaplanamaz. Tahmin ediciler, aynı veri üretim ve
örnekleme süreci çok kez yinelendiğinde gösterecekleri kuramsal veya benzetimsel
davranışla değerlendirilir.

\begin{table}[htbp]
\centering
\caption{Tahmin dilinin üç düzeyi.}
\begin{tabular}{@{}p{0.22\textwidth}p{0.28\textwidth}p{0.34\textwidth}@{}}
\toprule
Kavram & Gösterim & Yalnızlık araştırmasındaki karşılığı \\
\midrule
Parametre & $\mu$ veya $p$ & Evrenin bilinmeyen ortalaması veya risk oranı \\
Tahmin edici & $\bar X$ veya $\hat P$ & Örneklemden örnekleme değer alan hesaplama kuralı \\
Tahmin & $\bar x=43{,}6$, $\hat p=0{,}30$ & Gözlenen örneklemde hesaplanan sayı \\
\bottomrule
\end{tabular}
\end{table}

\begin{mistakebox}
Tahmin edici ile tahmini aynı şey gibi kullanmak kavramsal karışıklık yaratır.
$\bar X$ henüz örneklem seçilmeden önce rassal bir değişkendir; $\bar x$ ise
veri toplandıktan sonra elde edilen sabit sayıdır.
\end{mistakebox}

\section{Sık kullanılan nokta tahminleri}
\index{nokta tahmini!ortalama}
\index{nokta tahmini!oran}
\index{nokta tahmini!varyans}
\index{nokta tahmini!korelasyon}

\begin{table}[htbp]
\centering
\caption{Parametreler ve doğal örneklem tahminleri.}
\begin{tabular}{p{0.38\textwidth}p{0.36\textwidth}}
\toprule
Hedef parametre & Nokta tahmini \\
\midrule
Evren ortalaması $\mu$ & Örneklem ortalaması $\bar x$ \\
Evren oranı $p$ & Örneklem oranı $\hat p=x/n$ \\
Evren varyansı $\sigma^2$ & Örneklem varyansı $s^2$ \\
İki evren ortalaması farkı & $\bar x_1-\bar x_2$ \\
Korelasyon $\rho$ & Örneklem korelasyonu $r$ \\
\bottomrule
\end{tabular}
\end{table}

\subsection{Oran ve fark tahminleri}

Bir özelliği taşıyan $x$ gözlem $n$ kişilik örneklemde bulunmuşsa evren oranı
için
\[
\hat p=\frac{x}{n}
\]
kullanılır. İki grubun ortalama farkı için doğal tahmin
$\bar x_1-\bar x_2$, iki zaman ölçümündeki ortalama değişim için ise bireysel
fark puanlarının ortalamasıdır. Hangi yönde fark alındığı işaretin yorumunu
belirler.

Örneğin program grubunun ortalaması 38, karşılaştırma grubunun ortalaması 44
ise ``program eksi karşılaştırma'' farkı $38-44=-6$ puandır. Negatif işaret,
ilk yazılan grubun ortalamasının daha düşük olduğunu gösterir; etkinin iyi ya
da kötü olduğu sonucunu tek başına vermez.

Tek bir nokta tahmini belirsizliği göstermez. $\bar x=72$ sonucu, örneklem
hacmi 10 olduğunda ve 1000 olduğunda aynı görünür; oysa tahminlerin hassasiyeti
çok farklıdır. Bu nedenle nokta tahmini standart hata veya güven aralığıyla
birlikte verilmelidir.

\section{İyi bir tahmin edicide aranan özellikler}
\index{yansızlık}
\index{hassasiyet}
\index{etkinlik}
\index{tutarlılık}
\index{sağlamlık}

Tek bir özellik bütün durumlarda en iyi tahmin ediciyi belirlemez. Araştırmacı
aşağıdaki ölçütleri hedef ve veri yapısıyla birlikte değerlendirir:

\begin{description}
  \item[Yansızlık] Tekrarlı örneklemede tahminlerin merkezi hedef parametreye eşittir.
  \item[Hassasiyet] Tahmin edicinin örnekleme dağılımı dardır; standart hatası küçüktür.
  \item[Etkinlik] Benzer koşullardaki tahmin ediciler arasında daha düşük varyansa sahiptir.
  \item[Tutarlılık] Örneklem hacmi büyüdükçe hedef parametre çevresinde yoğunlaşır.
  \item[Sağlamlık] Aykırı değer veya küçük model sapmalarından aşırı etkilenmez.
\end{description}

Örneklem ortalaması normal veya yaklaşık simetrik veride etkili olabilir;
fakat güçlü aykırı değerlere duyarlıdır. Medyan daha sağlamdır, ancak simetrik
normal koşullarda ortalamadan daha değişken olabilir. Seçim ``hangi formül
daha gelişmiş?'' sorusuna değil, hedef parametre ve veri yapısına dayanır.

\section{Yanlılık}
\index{yanlılık!tahmin edici}

Bir tahmin edici için
\[
\text{Yanlılık}(\hat\theta)
=\E(\hat\theta)-\theta
\]
tanımlanır. Beklenen değeri hedef parametreye eşitse tahmin edici yansızdır.
Yansızlık uzun dönemli bir özelliktir; tek bir örneklemde tahminin parametreye
eşit olacağını söylemez.

\subsection{Tahmin edici yanlılığı ve veri yanlılığı}

Formülden kaynaklanan tahmin edici yanlılığı ile veri üretim sürecindeki
yanlılık ayrılmalıdır. Örneklem ortalaması evren ortalaması için kuramsal
olarak yansız olsa bile yalnız gönüllü öğrencilerden toplanan veride hedef
evren ortalamasını yanlı tahmin edebilir. Yansız bir formül, yanlı örnekleme
çerçevesini veya geçersiz ölçümü düzeltemez.

Örneklem varyansında paydanın $n-1$ olması da yansızlıkla ilişkilidir:
\[
s^2=\frac{1}{n-1}\sum_{i=1}^n(x_i-\bar x)^2.
\]
Sapmalar örneklem ortalaması çevresinde hesaplandığı için bir serbestlik
derecesi kullanılmış olur. $n$ paydasıyla hesaplanan değer, normal örnekleme
bağlamında evren varyansını sistematik olarak küçük tahmin eder.

\section{Varyans, standart hata ve ortalama karesel hata}
\index{tahmin edici!varyans}
\index{standart hata}
\index{ortalama karesel hata}
\index{MSE}

Tahmin edicinin örneklemden örnekleme değişkenliği varyansıyla ölçülür.
Standart hata bu varyansın kareköküdür. Yanlılık ile değişkenliği birleştiren
ortalama karesel hata
\[
MSE(\hat\theta)=\Var(\hat\theta)
+\text{Yanlılık}(\hat\theta)^2
\]
eşitliğini sağlar \cite{shao2003,wasserman2004}.

\begin{ornek}[Yanlılık--varyans ödünleşimi]
A tahmin edicisinin yanlılığı 0, varyansı 9; B tahmin edicisinin yanlılığı 2,
varyansı 2 olsun. Ortalama karesel hatalar
\[
MSE(A)=9+0^2=9,
\qquad
MSE(B)=2+2^2=6
\]
olur. B yanlı olmasına rağmen toplam karesel hata bakımından daha iyi sonuç
verebilir. Bu örnek yansızlığın tek başına yeterli bir seçim ölçütü olmadığını
gösterir.
\end{ornek}

\begin{sezgibox}
Hedef tahtasında atışların merkezi gerçek parametreyi temsil etsin. Atışların
merkezden sistematik kayması yanlılık, birbirinden uzaklığı varyanstır. İyi bir
tahmin edici hem hedefe yakın hem de tekrarlarda kararlı olmalıdır.
\end{sezgibox}

\section{Standart hata: Tahminin hassasiyeti}
\index{standart hata!ortalama}
\index{standart hata!oran}

Standart hata, tahmin edicinin tekrarlı örneklemedeki standart sapmasıdır.
Basit rastgele ve bağımsız örneklemede ortalama ve oran için yaklaşık
\[
SE(\bar X)=\frac{\sigma}{\sqrt n},
\qquad
SE(\hat p)=\sqrt{\frac{p(1-p)}{n}}
\]
olur. Bilinmeyen $sigma$ ve $p$ yerine uygulamada çoğunlukla $s$ ve
$\hat p$ kullanılır.

Yalnızlık örneğinde ortalamanın tahmini standart hatası
\[
\widehat{SE}(\bar X)=\frac{9{,}8}{\sqrt{120}}\approx0{,}89
\]
puandır. Risk oranının tahmini standart hatası ise
\[
\widehat{SE}(\hat p)
=\sqrt{\frac{0{,}30(0{,}70)}{120}}
\approx0{,}042
\]
olur. Bunlar bireysel puanların değişkenliğini değil, tahminlerin
hassasiyetini gösterir.

\begin{mistakebox}
Küçük standart hata, tahminin mutlaka doğru olduğu anlamına gelmez. Büyük ama
yanlı bir örneklem yanlış değeri çok hassas biçimde tahmin edebilir. Standart
hata örnekleme değişkenliğini ölçer; kapsam ve ölçüm yanlılığını içermez.
\end{mistakebox}

\begin{etkinlik}
Aynı ortalama ve standart sapmaya sahip $n=25$, $n=100$ ve $n=400$
örneklemlerinin standart hatalarını karşılaştırın. Hassasiyeti iki kat
artırmanın örneklem maliyetine etkisini tartışın.
\end{etkinlik}

\section{Tutarlılık}
\index{tutarlılık}

Örneklem hacmi büyüdükçe tahmin edici gerçek parametre çevresinde yoğunlaşıyorsa
tutarlı olarak adlandırılır. Örneklem ortalaması, sonlu varyans altında evren
ortalaması için tutarlıdır. Tutarlılık büyük örneklem özelliğidir ve küçük
örneklem performansını tek başına garanti etmez.

\section{Sağlam tahmin ve aykırı değer duyarlılığı}
\index{sağlam tahmin}
\index{medyan}
\index{kırpılmış ortalama}
\index{aykırı değer}

Ortalama her gözlemin büyüklüğünü kullandığı için uç değerlere duyarlıdır.
Medyan sıralamaya dayandığından daha sağlamdır. Örneğin
\[
20,22,23,24,25,26,27,80
\]
verisinin ortalaması 30,875; medyanı 24,5'tir. 80 değerinin kayıt hatası mı,
farklı bir alt gruba mı ait olduğu araştırılmadan hiçbir tahmin otomatik olarak
tercih edilmemelidir.

Kırpılmış ortalama, en küçük ve en büyük belirli orandaki gözlemleri dışarıda
bırakarak hesaplanır ve ortalama ile medyan arasında sağlam bir seçenek
sunabilir. Ancak dışlama oranı sonuç görüldükten sonra keyfî biçimde
seçilmemelidir.

\begin{researchbox}
PDR verilerinde uç puan araştırmanın en önemli katılımcısını temsil edebilir.
Sağlam tahmin kullanmak, uç gözlemi veri setinden silmekle aynı şey değildir.
Gözlem korunabilir; ortalama yanında medyan, IQR ve duyarlılık analizi
raporlanabilir.
\end{researchbox}

\section{Örnekleme tasarımı ve ağırlıklı tahmin}
\index{ağırlıklı tahmin}
\index{örnekleme ağırlığı}

Tabakalar farklı olasılıklarla seçilmişse basit örneklem ortalaması evren
bileşimini yansıtmayabilir. Evrenin yüzde 80'i A, yüzde 20'si B grubunda olsun.
Araştırmacı her gruptan eşit sayıda öğrenci seçmiş ve örneklem ortalamalarını
sırasıyla 40 ve 60 bulmuş olsun. Eşit örneklem sayılarıyla ağırlıksız ortalama
50'dir; evren paylarıyla ağırlıklı tahmin
\[
0{,}80(40)+0{,}20(60)=44
\]
olur. İki değer farklı hedefleri temsil eder. Alt grup karşılaştırmasında eşit
seçim yararlı olabilirken evren geneli tahminde ağırlık gerekebilir.

\begin{mistakebox}
Örneklemdeki yüzde ile evrendeki yüzde aynı değilse bütün satırlara eşit ağırlık
vermek evren ortalamasını çarpıtabilir. Ağırlık kullanımı örnekleme planı ve
hedef parametreyle birlikte açıklanmalıdır.
\end{mistakebox}

\section{Benzetimle yansızlık ve standart hata}

Aşağıdaki kod, ortalaması 50 ve standart sapması 10 olan bir evrenden 25
kişilik 5.000 örneklem seçer:

\begin{lstlisting}
import numpy as np

rng = np.random.default_rng(209)
tekrar = 5000
n = 25

ortalamalar = np.array([
    rng.normal(loc=50, scale=10, size=n).mean()
    for _ in range(tekrar)
])

print("Tahminlerin ortalamasi:", ortalamalar.mean())
print("Benzetim standart hatasi:", ortalamalar.std(ddof=1))
print("Kuramsal standart hata:", 10 / np.sqrt(n))
\end{lstlisting}

Tahminlerin ortalaması 50'ye, benzetimdeki standart sapma kuramsal
$10/\sqrt{25}=2$ değerine yaklaşır. Tek tek örneklem ortalamaları 50'ye eşit
olmak zorunda değildir; yansızlık tekrarlı örneklemin merkezine ilişkin bir
özelliktir.

\section{SPSS çıktısında nokta tahmini}

\begin{enumerate}
  \item \textbf{Analyze $>$ Descriptive Statistics $>$ Explore} menüsünü açın.
  \item Nicel değişkeni \textbf{Dependent List} alanına taşıyın.
  \item Ortalama, medyan, varyans, standart sapma ve ortalamanın standart hatasını inceleyin.
  \item Kategorik oran için \textbf{Analyze $>$ Descriptive Statistics $>$ Frequencies} yolunu kullanın.
  \item Geçerli ve eksik gözlem sayılarını tahminden önce kontrol edin.
\end{enumerate}

\begin{outputguidebox}
\textbf{Mean} gözlenen nokta tahminini, \textbf{Std. Deviation} bireysel
puanların yayılımını, \textbf{Std. Error} ise ortalama tahmininin hassasiyetini
gösterir. Bu üç sütun aynı kavram değildir. Geçerli $n$ değeri değişkenler
arasında farklıysa tahminlerin aynı katılımcı grubuna dayanmadığını unutmayın.
\end{outputguidebox}

\section{Bütünleştirici çözümlü örnek}

120 öğrencilik yalnızlık araştırmasında $\bar x=43{,}6$, $s=9{,}8$ ve yüksek
risk sayısı 36 bulunmuştur.

\begin{enumerate}
  \item Evren ortalaması için nokta tahmini 43,6 puandır.
  \item Evren risk oranı için nokta tahmini $36/120=0{,}30$, yani yüzde 30'dur.
  \item Ortalama için tahmini standart hata yaklaşık 0,89 puandır.
  \item Oran için tahmini standart hata yaklaşık 0,042, yani 4,2 yüzde puandır.
  \item Sonuç yalnız seçilen ve yanıtlayan öğrencilerin hedef evreni temsil ettiği ölçüde genellenebilir.
\end{enumerate}

Nokta tahminleri tek başına kesinlik izlenimi verebilir. Bir sonraki bölümde
standart hataların güven aralıklarına nasıl dönüştürüldüğü ele alınacaktır.

\section{Bilimsel raporlama örneği}

``Olasılıklı yöntemle seçilen 120 birinci sınıf öğrencisinin yalnızlık puanı
ortalaması 43,6'dır ($SS=9{,}8$, $SH=0{,}89$). Yüksek risk eşiğinin üzerinde
puan alan 36 öğrenci bulunmaktadır; örneklem oranı yüzde 30'dur
($SH=0{,}042$). Tahminler ölçme aracının tanımladığı puan ve araştırmanın hedef
evreniyle sınırlı olarak yorumlanmıştır.''

\section{Bir tahmini değerlendirirken}

\begin{enumerate}
  \item Hedef parametre açık mı?
  \item Örnekleme tasarımı hedef evrene uygun mu?
  \item Tahminin standart hatası veya güven aralığı verilmiş mi?
  \item Aykırı değerler ya da eksik veriler tahmini belirgin etkiliyor mu?
  \item Tahmin bağlama uygun birim ve hassasiyetle raporlanmış mı?
\end{enumerate}

\section{Literatür veri setiyle IBM SPSS laboratuvarı}
\index{Student Performance veri seti!nokta tahmini}
\index{SPSS!nokta tahmini}
\index{oran tahmini!ikili değişken ortalaması}
\index{standart hata!nokta tahmini}

Bu uygulamada \emph{Student Performance} matematik dosyasından iki hedef
büyüklük tahmin edilecektir \cite{cortez2008data,cortezsilva2008}:

\begin{enumerate}
  \item İlgili öğrenci topluluğundaki ortalama yıl sonu matematik notu,
  \item yıl sonu notu en az 10 olan öğrencilerin oranı.
\end{enumerate}

Örneklem ortalaması birinci parametre, 0--1 biçiminde kodlanan başarı
göstergesinin ortalaması ise ikinci parametre için nokta tahmini sağlar.
Başarı eşiği veri analizinden önce tanımlanmalı ve eğitim sistemindeki anlamı
kaynakla gerekçelendirilmelidir.

\begin{researchbox}
Veri setindeki 395 öğrenci iki okuldan gözlenmiştir. Aşağıdaki standart
hatalar, bu gözlemler uygun bir örneklem gibi ele alındığında oluşacak
örnekleme değişkenliğini betimler. Okulların veya öğrencilerin seçimindeki
olası kapsam ve seçim yanlılığını ölçmez; Portekiz'deki bütün öğrencilere
otomatik genelleme sağlamaz.
\end{researchbox}

\subsection{Araştırma görevleri}

\begin{enumerate}
  \item G3 ortalamasını, medyanını ve ortalamanın standart hatasını bulun.
  \item G3 puanı en az 10 olan öğrenciler için 0--1 başarı göstergesi üretin.
  \item Başarı göstergesinin ortalamasının neden örneklem oranına eşit olduğunu açıklayın.
  \item Ortalama ve oran tahminlerinin standart hatalarını ayırın.
  \item Tahminlerin hedef evrenini ve genelleme sınırını yazın.
\end{enumerate}

\subsection{SPSS syntax}

Eksik G3 değeri bulunması durumunda karşılaştırma işleminin eksik değeri yanlış
biçimde 0 olarak kodlamaması için başarı göstergesi önce sistem eksik olarak
oluşturulmuş, yalnız geçerli G3 değerlerinde hesaplanmıştır.

\begin{lstlisting}[language={}]
DATASET ACTIVATE StudentMath.

COMPUTE basarili=$SYSMIS.
IF NOT MISSING(G3) basarili=(G3>=10).
VARIABLE LABELS basarili 'Yıl sonu notu en az 10'.
VALUE LABELS basarili 0 'Başarısız' 1 'Başarılı'.
VARIABLE LEVEL basarili (NOMINAL).
FORMATS basarili (F1.0).
EXECUTE.

DESCRIPTIVES VARIABLES=G3 basarili
 /STATISTICS=MEAN SEMEAN STDDEV MIN MAX.

FREQUENCIES VARIABLES=G3
 /STATISTICS=MEDIAN
 /ORDER=ANALYSIS.

FREQUENCIES VARIABLES=basarili
 /ORDER=ANALYSIS.
\end{lstlisting}

Menüyle çalışırken G3 için \textbf{Analyze $>$ Descriptive Statistics $>$
Explore}; başarı oranı için \textbf{Transform $>$ Compute Variable} ve
ardından \textbf{Analyze $>$ Descriptive Statistics $>$ Frequencies} yolları
kullanılabilir. Menü seçimleri \textbf{Paste} ile syntax'a aktarılmalı ve
eşik tanımı analiz günlüğünde saklanmalıdır.

\subsection{SPSS çıktısının erişilebilir özeti}

\begin{table}[htbp]
\centering
\caption{G3 ortalaması ve başarı oranı için nokta tahminleri.}
\begin{tabular}{@{}p{0.34\textwidth}rrrrr@{}}
\toprule
Tahmin edilen büyüklük & Geçerli $n$ & Tahmin & SS & SH & Gözlenen aralık \\
\midrule
Ortalama G3 puanı & 395 & 10,415 & 4,581 & 0,231 & 0--20 \\
Başarılı öğrenci oranı & 395 & 0,671 & 0,470 & 0,024 & 0--1 \\
\bottomrule
\end{tabular}
\end{table}

\textbf{Frequencies} tablosunda 265 öğrenci başarılı, 130 öğrenci başarısız
olarak sınıflanır. Bu nedenle örneklem oranı
\[
\hat p=\frac{265}{395}=0{,}671
\]
veya yüzde 67,1'dir. G3'ün medyanı 11'dir; medyan, ortalamadan farklı bir
hedef büyüklüğün nokta tahminidir.

\subsection{Çıktının analizi}

G3 için \textbf{Mean}=10,415 örneklem ortalamasını, \textbf{Std.
Deviation}=4,581 öğrenciler arasındaki puan değişkenliğini ve \textbf{Std.
Error}=0,231 ortalama tahmininin örneklemden örnekleme değişkenliğini gösterir.
Standart hata
\[
SE(\bar X)=\frac{4{,}581}{\sqrt{395}}\approx0{,}231
\]
ile aynı sonucu verir. Tek bir öğrencinin puanının ortalamadan tipik uzaklığı
0,231 değil, yaklaşık 4,58 puandır.

Başarı göstergesi yalnız 0 ve 1 değerlerini aldığı için ortalaması 1'lerin
oranına eşittir. SPSS'in bu değişken için verdiği 0,470 standart sapma bireysel
başarı durumlarının değişkenliğini, yaklaşık 0,024 standart hata ise oran
tahmininin hassasiyetini gösterir:
\[
SE(\hat p)=\sqrt{\frac{0{,}671(1-0{,}671)}{395}}
\approx0{,}024.
\]

Ortalamanın standart hatasının küçük olması G3 ortalamasının tam olarak 10,415
olduğunu veya tahminin hedef evren için yanlı olmadığını kanıtlamaz. Standart
hata rassal örnekleme değişkenliğine ilişkindir; veri yalnız belirli okullardan
geldiğinde seçim mekanizması ayrıca tartışılmalıdır.

\begin{mistakebox}
Yüzde 67,1 başarı oranı ile 10,415 ortalama not aynı parametrenin iki farklı
gösterimi değildir. İlki önceden belirlenen eşiğe bağlı bir oran, ikincisi ise
notların bütün büyüklük bilgisini kullanan ortalamadır. Başarı eşiği 10'dan
12'ye değiştirilirse ortalama değişmez, fakat oran tahmini değişir.
\end{mistakebox}

\begin{outputguidebox}
Çıktıyı şu sırayla okuyun: (1) hedef parametre, (2) geçerli ve eksik $n$,
(3) nokta tahmini, (4) bireysel değişkenliği gösteren standart sapma,
(5) tahmin değişkenliğini gösteren standart hata, (6) kodlama ve eşik,
(7) hedef evren ve örnekleme sınırı. Küçük standart hatayı temsil gücünün
kanıtı olarak yorumlamayın.
\end{outputguidebox}

\subsection{Örnek bulgu paragrafı}

``İki okuldan gözlenen 395 öğrencinin yıl sonu matematik notu ortalama
10,42'dir ($SS=4{,}58$, $SH=0{,}23$); medyan 11'dir. Notu en az 10 olan 265
öğrenci bulunmaktadır ve örneklem başarı oranı yüzde 67,1'dir
($SH=2{,}4$ yüzde puan). Bu değerler ortalama not ve önceden tanımlanan başarı
oranı için nokta tahminleridir. Öğrencilerin seçimi bütün ülkeyi temsil eden
olasılıklı bir tasarıma dayanmadığından standart hatalar seçim yanlılığını
kapsamaz ve sonuçlar iki okul bağlamıyla sınırlı yorumlanmıştır.''

\section{Bölüm sonu çözümlü problemler}

\begin{enumerate}
  \item \textbf{Oran tahmini ve standart hata.} 200 öğrencinin 54'ü danışma hizmetini kullanmıştır. Nokta tahminini ve yaklaşık standart hatayı bulun.

  \textbf{Çözüm:} $\hat p=54/200=0{,}27$'dir. $SE(\hat p)=\sqrt{\hat p(1-\hat p)/n}=\sqrt{0{,}27(0{,}73)/200}\approx0{,}0314$ olur.

  \item \textbf{Yanlılık--varyans dengesi.} A tahmin edicisinin yanlılığı 1 ve varyansı 9; B'nin yanlılığı 2 ve varyansı 4 olsun. MSE değerlerini karşılaştırın.

  \textbf{Çözüm:} $MSE=\operatorname{Var}(\hat\theta)+\operatorname{Bias}(\hat\theta)^2$ olduğundan $MSE_A=9+1^2=10$, $MSE_B=4+2^2=8$'dir. B daha yanlı olmasına rağmen toplam karesel hata bakımından üstündür.

  \item \textbf{Ağırlıklı ortalama.} Evrenin yüzde 70'ini oluşturan A grubunun örneklem ortalaması 40, yüzde 30'unu oluşturan B grubunun ortalaması 55'tir. Evren ortalaması için tabaka ağırlıklı tahmini bulun.

  \textbf{Çözüm:} $\bar x_w=0{,}70(40)+0{,}30(55)=44{,}5$'tir. Basit grup ortalaması $(40+55)/2=47{,}5$, evren paylarını yanlışlıkla eşit kabul eder.
\end{enumerate}

\bolumkaynaklari{\cite{casella2002,hogg2019,shao2003,wasserman2004,cortez2008data,cortezsilva2008}}

\section*{Hatırla--Uygula--Yansıt}

\begin{enumerate}
\renewcommand{\labelenumi}{\textbf{\Alph{enumi}.}}
  \item \textbf{Hatırla:} Parametre, tahmin edici, tahmin ve tahmin hatasını aynı örnek üzerinde açıklayın.
  \item \textbf{Tanımla:} ``PDR öğrencilerinin tükenmişliği'' için açık bir hedef parametre yazın.
  \item \textbf{Hesapla:} 200 öğrencinin 46'sı bir hizmete başvurmuşsa oran tahminini ve tahmini standart hatayı bulun.
  \item \textbf{Karşılaştır:} Yansızlık, hassasiyet, etkinlik, tutarlılık ve sağlamlık kavramlarını ayırın.
  \item \textbf{MSE:} Yanlılığı 1, varyansı 8 olan A ile yanlılığı 2, varyansı 3 olan B tahmin edicisini MSE bakımından karşılaştırın.
  \item \textbf{Duyarlılık:} $12,13,14,15,16,50$ verisinde ortalama ve medyanı hesaplayıp uç değerin etkisini yorumlayın.
  \item \textbf{Ağırlıklandır:} Evren payları yüzde 70 ve 30, grup ortalamaları 45 ve 60 ise ağırlıklı evren ortalamasını bulun.
  \item \textbf{Eleştir:} ``Standart hata küçük olduğu için tahmin kesinlikle doğrudur'' ifadesini düzeltin.
  \item \textbf{Uygula:} Python benzetiminde $n=25$ yerine 100 kullanarak standart hatadaki değişimi inceleyin.
  \item \textbf{Raporla:} Bir ortalama ve bir oran tahminini örneklem hacmi, standart hata ve hedef evrenle birlikte yazın.
\end{enumerate}